\section{Proofs and additional details of the harmonic construction}
\label{app:harmonic-proofs}

\subsection{Reference coordinates and path preservation}

Let $\rho$ be a probability measure and let a globally defined smooth
flow generated by $b$ preserve $\rho$. For a proper smooth diffeomorphism
$F_t$, the pushforward vector field $w_t(F_tz)=DF_t(z)b(z)$ has zero
weighted divergence under $(F_t)_\#\rho$. Indeed,
\[
 \int\nabla\phi(F_tz)\cdot DF_t(z)b(z)\,\rho(\dd z)
 =\int\nabla(\phi\circ F_t)(z)\cdot b(z)\,\rho(\dd z)=0.
\]
Properness makes the pulled-back test compactly supported. The exact
flow of $w_t$ at frozen physical time is the conjugate of the flow of
$b$ by $F_t$. For $u_t(F_tz)=\partial_tF_tz$, material differentiation
along a fixed reference label yields
$\acc_t(F_tz)=\partial_{tt}F_tz$ and
$\jerk_t(F_tz)=\partial_{ttt}F_tz$.
Equation~\eqref{eq:reference-acceleration} therefore suffices exactly.

More generally, for a time-dependent scalar $\gamma(t)$, the field
$\partial_tF_t+\gamma(t)DF_tJz$ has trajectories
$F_t(R_{\int_0^t\gamma(s)\dd s}z)$, where $R_\theta z=e^{i\theta}z$.
For a radial reference law all these fields generate the same $p_t$.
In Eq.~\eqref{eq:harmonic-gauge}, $\gamma=\omega-a$.
Only the constant-$\omega$ family is used for the polynomial search.

\subsection{Global diffeomorphism, density, and regularity}

The differential of $f(z)=z^2e^{-|z|^2}$ in any unit real direction
$\xi$ obeys
\[
 |Df(z)\xi|
 \le 2|z|(1+|z|^2)e^{-|z|^2}\le4/e.
\]
The last scalar function is maximal at $|z|=1$, as follows by
differentiating it; its derivative has the sign of
$1+r^2-2r^4$. Hence
\begin{align*}
 |F_t(z)-F_t(y)|
 &\ge e^{at}(1-q)|z-y|,\\
 |F_t(z)-F_t(y)|
 &\le e^{at}(1+q)|z-y|.
\end{align*}
For every $x$ the map $z\mapsto A_t^{-1}(x-B_tf(z))$ is a contraction
of the complete space $\R^2$ with constant at most $q<1$. Its unique
fixed point proves surjectivity and gives the inverse. The differential
is nonsingular everywhere; the inverse function theorem gives a smooth
inverse. The singular values of $DF_t$ are between
$e^{at}(1-q)$ and $e^{at}(1+q)$. The determinant is positive, because
the perturbation may be continuously reduced to zero without crossing
a singular differential. The density formula in the main text follows
by change of variables.

Every derivative of $f$ is a polynomial times $e^{-|z|^2}$. It and its
products with any fixed power of $|z|$ are bounded. On a neighborhood
of the compact time interval, the strict inequality $q<1$ persists
after a small extension of the interval. The inverse maps have uniformly
bounded first derivative, and bounded higher derivatives of each fixed
order. Consequently each of $u_a$ and $\acc_{u_a}$ is a fixed linear
field plus a smooth bounded perturbation with bounded first derivatives.
They are globally Lipschitz, have at most linear growth, and the required
time derivatives exist on this neighborhood. Explicitly,
\[
 \acc_{u_a}(t,F_tz)=2ia^2A_tz+4ia^2B_tf(z),\qquad
 \jerk_{u_a}(t,F_tz)
 =((1+i)a)^3A_tz+((1+i)a\sqrt2)^3B_tf(z).
\]
Their $L^2(p_t)$ norms are uniformly finite. The same smoothness and
growth conclusions hold for each fixed $u_\omega$, and uniformly for
$\omega$ in compact sets. Gaussian moments give all the integrability
needed for the later Taylor and moment arguments. This proves the
regularity part of Theorem~\ref{thm:harmonic-construction}.

For the multi-mode extension, the sum of Lipschitz perturbation bounds
is strictly less than the leading linear singular value by hypothesis.
The same contraction argument applies. Angular equivariance gives
$\partial_\theta f_n=in f_n$, while
$\partial_{tt}(e^{(1+i)a\sqrt n t}f_n)=2ia^2n
e^{(1+i)a\sqrt n t}f_n$. Thus
$F_{tt}=2a^2F_\theta$ persists. Rapid decay of the $f_n$ and all their
derivatives, and finiteness of the sum, give the same regularity bounds.

\subsection{Why the example is neither affine-radial nor affine-product}

For a standard planar Gaussian $Z$, angular orthogonality gives
\begin{align}
 \E F_t(Z)&=0,\qquad \E F_t(Z)^2=0,\qquad \E F_t(Z)^3=0,
 \label{eq:harmonic-oddmoments}\\
 \E\bigl[F_t(Z)^2\overline{F_t(Z)}\bigr]
 &=A_t^2\overline B_t\,\E[|Z|^4e^{-|Z|^2}]
 =\frac8{27}A_t^2\overline B_t\ne0.\nonumber
\end{align}
Also $\E|F_t(Z)|^2=2|A_t|^2+(8/125)|B_t|^2$, so the real covariance
is a positive multiple of $I_2$. An affine image of a radial law is
centrally symmetric about its mean, which would make the last cubic
moment zero. Thus no such representation exists.

Suppose instead that the centered real vector $F_t(Z)$ were an
invertible affine image of independent scalar variables. Finite third
moments follow from those of $F_t(Z)$ and the inverse affine map.
After centering and standardizing the independent variables, isotropic
covariance forces the remaining mixing matrix to be a scalar multiple
of an orthogonal matrix. For independent centered real $U,V$,
\[
 \E(U+iV)^3=\E U^3-i\E V^3,\qquad
 \E[(U+iV)^2(U-iV)]=\E U^3+i\E V^3.
\]
The two quantities therefore have equal moduli. Orthogonal rotations,
reflections, and common scaling preserve this equality of moduli.
Equation~\eqref{eq:harmonic-oddmoments} violates it. This rules out an
affine-product representation, proving the final assertion of the
theorem. The small deformation is consequently not one of the earlier
radial, Gaussian, affine-radial, or independent-product constructions.

For Corollary~\ref{cor:gaussian-start}, the proposed bound is a uniform
upper bound on $(4/e)|B_t/A_t|$, because $e^{-at}\le1$ and
$|\sinh((1+i)a\sqrt2t)|\le\sqrt2\sinh(a\sqrt2T)$.
The diffeomorphism and regularity proofs therefore apply unchanged.
The second-derivative identity still holds, $B_0=0$, and $B_t\ne0$ for
$t>0$ since its real argument is positive. The moment argument above
continues to apply with this coefficient. This proves the corollary.

\subsection{Exact residual polynomial and correction recovery}

Along the $u_\omega$ trajectories the reference label rotates at rate
$\omega-a$. The material derivative on functions of $(t,z)$ is
$\partial_t+(\omega-a)\partial_\theta$. Applying it twice to each
$F_{n,t}$ gives the coefficient $\zeta_n^2$ in the main text.
Angular orthogonality makes different harmonics orthogonal in complex
$L^2(\rho)$, which is the real vector $L^2$ norm. Since
$|Z|^2$ has density $\tfrac12e^{-s/2}$ on $[0,\infty)$,
\[
 \E|Z|^2=2,\qquad
 \E[|Z|^4e^{-2|Z|^2}]
 =\frac12\int_0^\infty s^2e^{-5s/2}\dd s=\frac8{125}.
\]
This proves Eq.~\eqref{eq:harmonic-residual}. The first mode is
$\zeta_1=a+i\omega$. Its squared real part is $a^2-\omega^2$, so a
zero residual requires $\omega=\pm a$ and $c=2a\omega$.
If $\omega=-a$, the second mode has real part of its squared
coefficient $2a^2-a^2(\sqrt2-4)^2\ne0$. If $\omega=a$, both mode
coefficients equal $inc$ for $c=2a^2$. As all $M_n$ are positive,
these alternatives prove uniqueness of the zero.

For fixed $\omega$, differentiate the quadratic objective in $c$ to
obtain $c_*$. The imaginary part of $\zeta_n^2$ is affine in $\omega$,
so eliminating $c$ changes only terms of degree at most two in
$\omega$. The real squared coefficients contribute the positive
quartic leading term $\sum n^4M_n$. Coercivity and differentiation
therefore give the finite global minimization procedure. The same
calculation applies to time integrals, because the coefficients
$\zeta_n$ are time-independent. This proves
Proposition~\ref{prop:harmonic-search}.

The upper certificate controls the canonical projected residual by
orthogonal projection. It need not equal that residual away from the
recovered zero, because a single solenoidal basis field need not be the
best approximation among all solenoidal fields. The proof certifies
the exact algebraic solution $\omega=a,c=2a^2$; machine residuals near
roundoff describe a numerical recovery of this solution.

\subsection{A finite modified-flow argument for the explicit coupling}

Let $u=u_a$ and $\acc=\acc_u$. For each $0<h\le h_0$, define
$u_h=u-(h/2)\acc$. The curve
\[
 Y_h(t,z)=F_t(e^{-ia^2ht}z)
\]
solves $\dot Y_h=u_h(t,Y_h)$ and starts from $F_0(z)$.
For $Z\sim\rho$, its law at each $t$ is exactly $p_t$.
On a fixed bounded range of $h$, $u_h$ and its required derivatives
satisfy uniform linear-growth and spatial Lipschitz bounds by the
previous regularity calculation. Its material acceleration is
\[
 \partial_tu_h+(D u_h)u_h
 =\acc-\frac h2\{\jerk+(Du)\acc\}
       +\frac{h^2}4(D\acc)\acc,
\]
All remainder terms have an $L^2$ bound proportional to
$1+|z|$ along $Y_h$, uniformly in $h$ and $t$; the third derivatives
needed for the Taylor remainder have the same bound. Taylor expansion
over one step gives
\begin{align*}
 Y_h(t+h,z)
 &=Y_h(t,z)+h[u(t,Y_h)-(h/2)\acc(t,Y_h)]\\
 &\quad+\frac{h^2}2\acc(t,Y_h)
       +O_{L^2(\rho)}(h^3)\\
 &=Y_h(t,z)+h u(t,Y_h(t,z))+O_{L^2(\rho)}(h^3).
\end{align*}
Minkowski's inequality and Euler stability imply
$e_{k+1}\le(1+Lh)e_k+Ch^3$, where
$e_k=\|\widehat X_k-Y_h(t_k,Z)\|_{L^2}$ and $e_0=0$.
Summing yields $\max_ke_k\le CT e^{LT}h^2$. Because rotations preserve
$\rho$, the pairing at $t=T$ proves
Eq.~\eqref{eq:harmonic-coupling}.

Finally, Taylor expansion in the rotation angle gives in $L^2(\rho)$
\[
 F_T(e^{-ia^2Th}Z)
 =F_T(Z)-a^2Th\,DF_T(Z)JZ+O_{L^2}(h^2).
\]
The preceding $O(h^2)$ coupling error therefore shows
$h^{-1}(\widehat X_N-F_T(Z))\to-a^2TDF_T(Z)JZ$ in $L^2$.
Taking norms proves Eq.~\eqref{eq:harmonic-particle-limit}. Its constant
is positive since $DF_T$ is nonsingular and $\rho(Z=0)=0$.
For a uniform grid the same argument works at every grid point with
$T$ replaced by $t_k$. The law-level $O(H^2)$ assertion for arbitrary
grids is supplied separately by Theorem~\ref{thm:distribution_order}.
